\documentclass{article}
\usepackage[utf8]{inputenc}
\usepackage[T1]{fontenc}
\usepackage{amsmath}
\usepackage{amssymb}
\usepackage{amsthm}
\usepackage{array}
\usepackage[portuguese]{babel}
\usepackage{booktabs}

% --- Configuração dos Ambientes ---
\newtheoremstyle{definicao_estilo}
  {\topsep} % Espaço acima
  {\topsep} % Espaço abaixo
  {}% Estilo da fonte do corpo
  {}% Recuo
  {\bfseries}   % Estilo da fonte do cabeçalho
  {.}   % Pontuação após o cabeçalho
  {.5em}% Espaço após o cabeçalho
  {}% Especificação do cabeçalho (padrão)


% Cria o ambiente "Definição"
\theoremstyle{definicao_estilo}
\newtheorem{definicao}{Definição}

 % Define um ambiente de propriedade, prova, resultado
\newtheorem{propriedade}{Propriedade}
\newtheorem{prova}{Prova}
\newtheorem{resultado}{Resultado}


% Inicio do Documento
\title{Criptografia: Conceitos Fundamentais e Curvas Elípticas}
\author{Vinicius Tavares}
\date{\today}

\begin{document}

\maketitle

% ------------------------------------------------------------------------------------------
\section{Criptografia}

Criptografia é a prática e o estudo de técnicas para comunicação segura na presença de comportamento adversário. De forma mais geral, a criptografia trata da construção e análise de protocolos que impedem terceiros ou o público de ler mensagens privadas.

% ------------------------------------------------------------------------------------------
% Exemplo modificado do livro "Guide to Elliptic Curve Cryptograph"
\par\vspace{1em}

\textbf{\textit{Exemplo Clássico: O Modelo do Adversário (Alice, Bob e Eva)}}

\par\vspace{0.5em}

\textbf{\textit{O Cenário}}
\begin{itemize}
\item \textbf{Alice ($A$)} e \textbf{Bob ($B$)} são agentes secretos que precisam coordenar uma operação.
\item \textbf{Eva ($E$)} é a oficial de inteligência inimiga. Sua única tarefa é interceptar e decifrar todas as comunicações.
\item \textbf{O Canal Não Seguro:} Alice e Bob usam um rádio de transmissão pública, monitorado 24 horas por dia por Eva.
\end{itemize}

\par\vspace{0.5em}

\textbf{\textit{O Segredo Compartilhado (A Chave Simétrica)}}
Antes da missão, Alice e Bob concordaram em usar uma \textbf{chave secreta} ($K$) que Eva não conhece, baseada em um sistema de coordenadas para um livro específico.
\begin{itemize}
\item \textbf{A Chave ($K$):} O livro \textit{Os Lusíadas} de Camões e a instrução para decodificar cada bloco de \textbf{7 dígitos} ($C$) da seguinte forma: os \textbf{3 primeiros} dígitos representam a página do livro, os \textbf{2 seguintes} a linha, e os \textbf{2 últimos} a palavra na linha.
\item \textbf{A Mensagem Original ($M$):} "Ataque ao nascer do sol." (reduzida para "Ataque Nascer Sol")
\end{itemize}

\par\vspace{0.5em}

\textbf{\textit{O Modelo do Adversário (Eva)}}
O \textbf{Modelo do Adversário} define o poder de Eva e as suposições de segurança do sistema.
\begin{center}
\begin{tabular}{|>{\bfseries}l|p{4.5cm}|p{4.5cm}|}
\hline
\textbf{Capacidade de Eva} & \textbf{Descrição} \\
\hline
\textbf{Poder de Interceção} & Eva \textbf{intercepta} toda a comunicação (os sinais de rádio), recebendo sempre a \textbf{mensagem cifrada} ($C$). \\
\hline
\textbf{Conhecimento do Algoritmo} & Eva sabe que está sendo usada uma cifra de código. Em criptografia, assumimos que o adversário \textbf{conhece o algoritmo}. \\
\hline
\textbf{Objetivo} & O objetivo de Eva é obter a \textbf{Mensagem Original} ($M$). \\
\hline
\textbf{Limitação} & Eva \textbf{NÃO} conhece a Chave Secreta ($K$). \\
\hline
\end{tabular}
\end{center}

\par\vspace{0.5em}

\textbf{\textit{A Troca de Mensagens (O Algoritmo)}}

\par\vspace{0.2em}
\textbf{\textit{1. A Ação de Alice (Cifragem)}}
Alice usa a chave ($K$) para transformar a mensagem ($M$) em uma série de blocos de sete dígitos ($C$), onde cada bloco representa uma palavra (coordenada).
\begin{itemize}
\item $M \to C$ (Ex: $5112703 \ 1590308 \ 1590310$)
\item Alice envia os blocos de sete dígitos pelo rádio público.
\end{itemize}

\par\vspace{0.2em}
\textbf{\textit{2. A Interceção de Eva (Eavesdropping)}}
Eva ouve o rádio e anota os números: $C = 0040507 \ 1280914 \ 0015003$. Ela sabe que é uma mensagem secreta, mas os blocos de 7 dígitos não têm sentido para ela.

\par\vspace{0.2em}
\textbf{\textit{3. A Ação de Bob (Decifragem)}}
Bob recebe os mesmos números ($C$).
\begin{itemize}
\item Bob aplica a Chave Secreta ($K$) no processo inverso.
\item Ele usa o livro \textit{Os Lusíadas} e a regra de coordenadas da chave para localizar o vocábulo correspondente a cada bloco.
\item $C \to M$ (Ex: $5112703 \ 1590308 \ 1590310$ se transforma em "Ataque Nascer Sol.")
\end{itemize}

\par\vspace{0.5em}

\textbf{\textit{A Falha de Eva}}
Eva não consegue ler a mensagem antes que a operação seja concluída, pois não tem a \textbf{Chave Secreta ($K$)}. O tempo necessário para tentar adivinhar a chave (a \textbf{complexidade computacional}) é inviável.

\par
\textbf{Conclusão:} O sistema é seguro \textbf{sob o Modelo de Adversário} onde Eva é \textbf{limitada por não ter a chave secreta} e pela dificuldade matemática de encontrá-la.

% ------------------------------------------------------------------------------------------
\section{Objetivos de Segurança}


Com esse exemplo, podemos ilustrar alguns objetivos fundamentais que uma comunicação segura busca atingir:

\begin{enumerate}
\item \textbf{Confidencialidade:} Manter os dados em segredo de todos, exceto daqueles autorizados a vê-los. Mensagens enviadas de $A$ para $B$ não devem ser legíveis por $E$.
\item \textbf{Integridade dos Dados:} Garantir que os dados não foram alterados por meios não autorizados. $B$ deve ser capaz de detectar quando os dados enviados por $A$ foram modificados por $E$.
\item \textbf{Autenticação da Origem dos Dados:} Corroborar a fonte dos dados. $B$ deve ser capaz de verificar que os dados supostamente enviados por $A$ de fato se originaram em $A$.
\item \textbf{Autenticação de Entidade:} Corroborar a identidade de uma entidade. $B$ deve estar convencido da identidade da outra entidade comunicante.
\item \textbf{Não-Repúdio:} Prevenir que uma entidade negue compromissos ou ações anteriores. $A$ não pode negar ter enviado a mensagem para $B$, e $B$ pode provar isso a um terceiro neutro.
\end{enumerate}
Algumas aplicações podem ter outros objetivos de segurança, como \textit{anonimato} ou \textit{controle de acesso}.

% ------------------------------------------------------------------------------------------
\section{Criptografia de Chave Simétrica vs. Assimétrica}
Sistemas de criptografia podem ser majoritariamente divididos em dois tipos: Em \textit{esquemas de chave simétrica} ou \textit{esquemas de chave assimétrica (ou chave pública)}.

\subsection{Criptografia de Chave Simétrica}
Os agentes comunicantes utilizam uma \textbf{chave secreta única} ($K$) para os processos de cifragem e decifragem.
\begin{itemize}
\item \textbf{Cifragem:} $C = E_K(M)$
\item \textbf{Decifragem:} $M = D_K(C)$
\end{itemize}
É extremamente \textbf{eficiente e rápida}. A principal desvantagem é o \textbf{problema da distribuição de chaves}, ou seja, como $A$ e $B$ podem concordar secretamente sobre $K$ em um canal inseguro.

\subsection{Criptografia de Chave Assimétrica (ou Chave Pública)}
Os agentes comunicantes utilizam um \textbf{par de chaves} interligadas, mas distintas: uma \textbf{Chave Pública} ($K_{pub}$), distribuída livremente, e uma \textbf{Chave Privada} ($K_{priv}$), mantida em segredo.
\begin{itemize}
\item \textbf{Cifragem (para B):} $C = E_{K_{pub\_B}}(M)$
\item \textbf{Decifragem (por B):} $M = D_{K_{priv\_B}}(C)$
\end{itemize}
Este método resolve o problema da distribuição de chaves e é crucial para a autenticação (\textbf{assinatura digital}). A chave pública é como a fenda de uma caixa de correio (qualquer um deposita), e a chave privada é a única chave que abre o compartimento.

% ------------------------------------------------------------------------------------------
\section{Sistemas RSA}
O RSA (Rivest, Shamir e Adleman), proposto em 1977, é um dos primeiros e mais importantes sistemas de criptografia de chave pública.

\subsection{Funcionamento do RSA}
A segurança do sistema RSA reside na dificuldade computacional de fatorar grandes números primos. O RSA utiliza a aritmética modular, onde todas as operações são realizadas no anel de inteiros $\mathbb{Z}_N$, ou seja, o resultado é sempre o resto da divisão por $N$.

\begin{itemize}
\item \textbf{O Módulo ($N$):} O número $N$ é o \textbf{Módulo RSA}, o resultado do produto de dois números primos muito grandes ($p$ e $q$). Ele define o tamanho do "universo" matemático no qual a cifra opera. O tamanho da chave pública e privada é ditado pelo tamanho de $N$.
\item \textbf{A Chave Pública ($e, N$):} O par $(e, N)$ é conhecido por todos e utilizado para cifrar a mensagem.
\item \textbf{A Chave Privada ($d$):} O expoente $d$ é mantido em segredo e utilizado para decifrar.
\end{itemize}

O processo de cifragem e decifragem é baseado nas seguintes operações de exponenciação modular:

\begin{enumerate}
\item \textbf{Preparação da Mensagem ($M$):} A mensagem original é convertida em um inteiro $M$ tal que $0 \le M < N$. Caso $M$ seja maior que $N$, é dividida em blocos menores.
\item \textbf{Cifragem (Processo Público):} Para cifrar a mensagem $M$, o emissor (Alice) utiliza a Chave Pública $(e, N)$ do receptor e calcula a mensagem cifrada $C$ pela seguinte fórmula:
$$ C \equiv M^e \pmod N $$
\item \textbf{Decifragem (Processo Privado):} Para decifrar a mensagem cifrada $C$, o receptor (Bob) utiliza sua Chave Privada $d$ e o Módulo $N$ para recuperar a mensagem original $M$:
$$ M \equiv C^d \pmod N $$
\end{enumerate}

\subsubsection*{O Princípio de Correção do RSA e a Função Alçapão}
A decifragem funciona porque o expoente privado $d$ é o inverso multiplicativo do expoente público $e$ no contexto da função totiente de Euler, $\phi(N) = (p-1)(q-1)$. A relação fundamental que garante que $C^d$ retorna $M$ é o \textbf{Teorema de Euler} generalizado:

\begin{propriedade}
Se $N=pq$ é o produto de dois primos e $ed \equiv 1 \pmod{\phi(N)}$, então para qualquer mensagem $M$:
$$ M^{ed} \equiv M \pmod N $$
\end{propriedade}

A segurança do RSA baseia-se no conceito de \textbf{Função Alçapão} (conceito que será melhor definido adiante). A operação de cifragem ($M \to C$) é fácil de executar. A operação inversa (decifragem $C \to M$) é computacionalmente intratável (muito difícil) sem o conhecimento do "alçapão" ($d$). O alçapão $d$ só pode ser calculado se os fatores primos $p$ e $q$ de $N$ forem conhecidos, o que equivale ao problema da fatoração de grandes números.

\subsection{Geração de Chave RSA}
Um par de chaves RSA consiste em uma chave pública $(e,N)$ e uma chave privada $d$. O processo de geração é baseado na dificuldade de fatorar grandes números.

\subsubsection{Passos para Geração da Chave}
A chave pública é o par de inteiros $(e, N)$.
\begin{itemize}
\item O \textbf{Módulo RSA} ($N$) é o produto de dois números primos ($p$ e $q$) grandes e aleatórios: $N = pq$.
\item O \textbf{Expoente de Cifragem} ($e$) é um inteiro que satisfaz $1 < e < \phi(N)$ e $\text{mdc}(e, \phi(N)) = 1$.
\item O \textbf{Expoente de Decifragem} ($d$), a chave privada, é um inteiro que satisfaz $1 < d < \phi(N)$ e a congruência:
$$ ed \equiv 1 \pmod{\phi (N)} $$
onde $\phi(N)$ é a função totiente de Euler: $\phi(N) = (p-1)(q-1)$.
\end{itemize}

\subsubsection{Pseudo Código de Geração}
\begin{verbatim}
Input: parâmetro de segurança l.
Output: Chave pública RSA (N, e) e chave privada d.
  1. Seleciona aleatoriamente dois primos p e q.
  2. Calcula N = pq e phi = (p-1)(q-1).
  3. Seleciona um inteiro e arbitrário com 1 < e < phi e mdc(e, phi) = 1.
  4. Calcula o inteiro d que satisfaz 1 < d < phi e ed = 1 (mod phi).
  5. Retorna (N, e, d).
\end{verbatim}

\section{Sistemas De Logaritmo Discreto}
A base de segurança para muitos sistemas criptográficos assimétricos, incluindo o Diffie-Hellman e os sistemas baseados em curvas elípticas, é o \textbf{Problema do Logaritmo Discreto} (DLP - Discrete Logarithm Problem).

\subsection{O Problema do Logaritmo Discreto (DLP)}
Seja $G$ um grupo cíclico de ordem $n$, e seja $g$ um gerador de $G$. Para um dado $x \in \mathbb{Z}_n$, calcular o elemento do grupo $y = g^x$ é computacionalmente fácil (problema da exponenciação discreta). No entanto, dado $g$ e $y$, encontrar o expoente $x$ tal que $y = g^x$ é considerado um problema computacionalmente intratável em grupos grandes o suficiente.

\begin{definicao}
O \textbf{Logaritmo Discreto} $x$, dado os elementos do grupo $g$ e $y$, é o expoente único $x \in \{0, 1, \dots, n-1\}$ tal que $g^x = y$.
\end{definicao}

A dificuldade de resolver o DLP em grupos específicos é o que confere segurança a estes sistemas.

\subsection{Função Alçapão (Trapdoor Function)}
O Logaritmo Discreto, juntamente com a Fatoração de Inteiros (base do RSA), constitui um tipo especial de função matemática central para a criptografia de chave pública: a \textbf{Função Alçapão}.

\begin{definicao}
Uma \textbf{Função Alçapão} (ou \emph{Trapdoor Function}) é uma função que é fácil de calcular em uma direção, mas extremamente difícil de inverter sem possuir informações secretas específicas (o "alçapão").
\end{definicao}

No contexto do DLP, a função $f(x) = g^x$ é a direção "fácil" (cifragem), enquanto a informação secreta (o alçapão) é a chave privada $x$ (o logaritmo discreto) que permite a inversão (decifragem) eficiente. Se o alçapão é conhecido, a dificuldade computacional é superada.

\section{Sistema de Curvas Elípticas}

A Criptografia de Curvas Elípticas (ECC - Elliptic Curve Cryptography) é uma abordagem moderna para criptografia de chave pública baseada na estrutura algébrica de curvas elípticas sobre campos finitos. Sua segurança deriva da dificuldade de resolver o \textbf{Problema do Logaritmo Discreto em Curvas Elípticas} (ECDLP), que é significativamente mais difícil de resolver do que o DLP tradicional em grupos finitos ou o problema da fatoração (base do RSA) para o mesmo tamanho de chave. Isso permite que a ECC ofereça o mesmo nível de segurança que outros esquemas com chaves muito menores (e.g., uma chave ECC de 256 bits é comparável em segurança a uma chave RSA de 3072 bits), resultando em maior eficiência e menor consumo de recursos.

\subsection{Grupo de Curvas Elípticas}
\textbf{\textit{Definição (Grupo De Curvas Elípticas)}}

Seja $p$ um número primo, e seja $\mathbb{F}_p$ o campo dos inteiros módulo $p$. Uma \emph{curva elíptica $E$ sobre $\mathbb{F}_p$} é definida por uma equação da forma
\begin{equation}
y^2 = x^3 + ax + b,
\tag{1.4}\label{eq:curva_eliptica}
\end{equation}
onde $a, b \in \mathbb{F}_p$ satisfazem $4a^3 + 27b^2 \not\equiv 0 \pmod p$. Um par $(x, y)$, onde $x, y \in \mathbb{F}_p$, é um \emph{ponto} na curva se $(x, y)$ satisfaz a equação (\ref{eq:curva_eliptica}). O \emph{ponto no infinito}, denotado por $\infty$, também é considerado um ponto na curva. O conjunto de todos os pontos em $E$ é denotado por $E(\mathbb{F}_p)$.

\subsection{A Soma de Pontos em $E(\mathbb{F}_p)$}

O conjunto de pontos em uma curva elíptica $E$ sobre $\mathbb{F}_p$, denotado por $E(\mathbb{F}_p)$, forma um grupo abeliano sob uma operação binária chamada "adição de pontos" (ou soma). O elemento neutro (ou identidade) é o \textbf{Ponto no Infinito} ($\infty$), que atua como o \textbf{zero} desta soma.

\subsubsection{O Ponto no Infinito ($\infty$) - Elemento Neutro}

Na geometria da curva elíptica, o $\mathbf{\infty}$ é crucial para completar a estrutura do grupo.

\begin{itemize}
\item \textbf{Elemento Neutro:} Assim como zero em números reais ($5+0=5$), a soma de qualquer ponto com o Ponto no Infinito é o próprio ponto:
$$ P_1 + \infty = P_1 $$
\item \textbf{Inverso (Negativo):} O inverso de um ponto $P_1 = (x_1, y_1)$ é o ponto $\mathbf{-P_1} = (x_1, -y_1)$, que é a sua reflexão sobre o eixo $x$.
\item \textbf{Soma de Inversos:} Quando a reta que une um ponto $P_1$ ao seu inverso $-P_1$ é traçada, ela é vertical. A definição geométrica da soma afirma que uma reta vertical intersecta a curva em $P_1$, $-P_1$ e, finalmente, no Ponto no Infinito ($\infty$).
\item \textbf{Regra:} A soma de um ponto com seu inverso resulta no elemento neutro:
$$ P_1 + (-P_1) = \infty $$
\end{itemize}

\subsubsection*{Soma de Dois Pontos Distintos ($P_1 \neq P_2$)}
Se $P_1$ e $P_2$ são distintos e $x_1 \neq x_2$, a soma $R = P_1 + P_2 = (x_3, y_3)$ é calculada em três passos:

\begin{enumerate}
\item \textbf{Cálculo do Coeficiente Angular ($\lambda$):}
$$ \lambda = \frac{y_2 - y_1}{x_2 - x_1} \pmod p $$
\item \textbf{Coordenada $x$ do Resultado ($x_3$):}
$$ x_3 = \lambda^2 - x_1 - x_2 \pmod p $$
\item \textbf{Coordenada $y$ do Resultado ($y_3$):}
$$ y_3 = \lambda(x_1 - x_3) - y_1 \pmod p $$
\end{enumerate}

Em termos simples, a soma de dois pontos $P$ e $Q$ em uma curva elíptica não é uma soma numérica tradicional, mas sim uma \textbf{operação geométrica} que garante que o resultado seja sempre um novo ponto na própria curva, transformando o conjunto de pontos em um grupo matemático.

Geometricamente, somar dois pontos $P_1$ e $P_2$ (com $P_1 \neq P_2$) pode ser definido como: 
\begin{enumerate}
\item Traçar uma reta que passa por $P_1$ e $P_2$.
\item Encontrar o terceiro ponto $R'$ no qual a reta intercepta a curva.
\item Refletir $R'$ com respeito ao eixo $x$ e encontrar $R$.
\end{enumerate}

\subsubsection*{Dobra de um Ponto ($P_1 = P_2$)}
Se desejamos calcular $2P_1 = P_1 + P_1$, a fórmula do coeficiente angular ($\lambda$) muda para a inclinação da reta tangente ao ponto $P_1$:

\begin{enumerate}
\item \textbf{Cálculo do Coeficiente Angular ($\lambda$):}
Lembrando que a curva é $y^2 = x^3 + ax + b$, a inclinação é dada pela derivada implícita:
$$ \lambda = \frac{3x_1^2 + a}{2y_1} \pmod p $$
\item \textbf{Coordenada $x$ do Resultado ($x_3$):}
$$ x_3 = \lambda^2 - 2x_1 \pmod p $$
\item \textbf{Coordenada $y$ do Resultado ($y_3$):}
$$ y_3 = \lambda(x_1 - x_3) - y_1 \pmod p $$
\end{enumerate}

Novamente, dobrar um ponto $P$ pode ser definido de maneira geométrica como:

\begin{enumerate}
\item Traçar a reta tangente à curva exatamente no ponto $P$.
\item Encontrar o segundo ponto $R'$ no qual a reta tangente intercepta a curva.
\item Refletir $R'$ com respeito ao eixo $x$ e encontrar $R = 2P$.
\end{enumerate}

\subsection{Multiplicação Escalar e o ECDLP}

A operação fundamental na ECC é a \textbf{Multiplicação Escalar}, que é essencialmente a soma repetida de um ponto. O ECDLP é baseado na dificuldade de inverter esta operação.

\subsubsection{Multiplicação Escalar}
A multiplicação escalar de um ponto $P$ por um escalar $k$ (um inteiro) é definida como a soma de $P$ consigo mesmo $k$ vezes:
$$ Q = kP = \underbrace{P + P + \dots + P}_{k \text{ vezes}} $$
Esta operação é realizada de forma eficiente através de algoritmos como a \textbf{Adição e Dobra Binária} (ou algoritmo de exponenciação modular adaptado).

\subsubsection{O Problema do Logaritmo Discreto em Curvas Elípticas (ECDLP)}
A segurança da ECC repousa sobre a dificuldade computacional de inverter a multiplicação escalar.

\begin{definicao}
O \textbf{Problema do Logaritmo Discreto em Curvas Elípticas (ECDLP)} é o seguinte: Dado um ponto base $G$ e o resultado da multiplicação escalar $Q = kG$, encontrar o escalar (inteiro) $k$ é computacionalmente intratável se a curva e os parâmetros forem escolhidos adequadamente.
\end{definicao}

No contexto da criptografia de chave pública:
\begin{itemize}
\item $k$ é a \textbf{chave privada} (o alçapão).
\item $G$ é o \textbf{ponto base} (parâmetro público).
\item $Q$ é a \textbf{chave pública} (resultado da operação).
\end{itemize}

\subsection{Uso Criptográfico: Geração e Troca de Chaves}
Para utilizar a ECC, todos os participantes devem concordar em um conjunto de parâmetros globais.

\subsubsection{Parâmetros de Domínio da Curva Elíptica}
Um esquema ECC requer a definição de um conjunto de seis parâmetros de domínio $\mathcal{D} = (p, a, b, G, n, h)$, que devem ser públicos e verificados:
\begin{itemize}
\item $p$: O módulo primo que define o campo finito $\mathbb{F}_p$.
\item $a, b$: Coeficientes que definem a equação da curva $E: y^2 = x^3 + ax + b$.
\item $G$: O \textbf{ponto base} (ou gerador) na curva. Todos os cálculos de chaves públicas partem deste ponto.
\item $n$: A \textbf{ordem} do ponto $G$ (o menor inteiro positivo tal que $nG = \infty$). Este $n$ deve ser um número primo grande.
\item $h$: O cofator, $h = |E(\mathbb{F}_p)|/n$, onde $|E(\mathbb{F}_p)|$ é o número total de pontos na curva.
\end{itemize}

\subsubsection{Geração de Chaves ECC}
A geração de chaves é um processo de duas etapas realizado por cada usuário (e.g., Alice).

\begin{enumerate}
\item \textbf{Seleção da Chave Privada ($k_{priv}$):} Alice seleciona um inteiro $k_{priv}$ (seu escalar secreto) aleatório no intervalo $[1, n-1]$.
\item \textbf{Cálculo da Chave Pública ($P_{pub}$):} Alice calcula o ponto $P_{pub}$ na curva usando o ponto base $G$ e sua chave privada:
$$ P_{pub} = k_{priv} \cdot G $$
O ponto $P_{pub}$ é sua chave pública.
\end{enumerate}

\subsubsection{Pseudo Código para Geração de Chave ECC}
\begin{verbatim}
Entrada: Parâmetros de Domínio (G, n)
Saída: Chave Privada (k_priv), Chave Pública (P_pub)

1. k_priv = Seleciona_Aleatorio(1, n-1)
2. P_pub = Multiplicacao_Escalar(k_priv, G)

3. Retorna (k_priv, P_pub)
\end{verbatim}

\subsection{Exemplo de Aplicação: Esquema de Troca de Chaves Diffie-Hellman de Curva Elíptica (ECDH)}
O ECDH é um protocolo fundamental usado para estabelecer um segredo compartilhado (chave de sessão) entre duas partes em um canal inseguro.

\subsubsection{Processo ECDH (Alice e Bob)}
Assumindo que Alice e Bob concordam nos parâmetros de domínio $(G, n)$, o processo é o seguinte:

\begin{center}
\begin{tabular}{|>{\bfseries}p{1.8cm}|p{4.5cm}|p{4.5cm}|}
\hline
\textbf{Etapa} & \textbf{Alice (Chave Privada $k_A$)} & \textbf{Bob (Chave Privada $k_B$)} \\
\hline
\textbf{1. Chaves Privadas} & Seleciona $k_A \in [1, n-1]$ & Seleciona $k_B \in [1, n-1]$ \\
\hline
\textbf{2. Chaves Públicas} & Calcula $P_A = k_A \cdot G$ & Calcula $P_B = k_B \cdot G$ \\
\hline
\textbf{3. Troca} & Envia $P_A$ para Bob & Envia $P_B$ para Alice \\
\hline
\textbf{4. Segredo Compartilhado} & Calcula $S = k_A \cdot P_B$ & Calcula $S = k_B \cdot P_A$ \\
\hline
\end{tabular}
\end{center}

\subsubsection{Prova de Igualdade do Segredo}
O segredo $S$ calculado por ambos é idêntico devido à propriedade associativa da multiplicação escalar:
$$ S_{\text{Alice}} = k_A \cdot P_B = k_A \cdot (k_B \cdot G) $$
$$ S_{\text{Bob}} = k_B \cdot P_A = k_B \cdot (k_A \cdot G) $$
Portanto, $S_{\text{Alice}} = S_{\text{Bob}} = (k_A k_B) \cdot G$.

\subsubsection{Segurança do ECDH}
Um adversário (Eva) intercepta as chaves públicas $P_A$ e $P_B$ e o ponto base $G$. Para calcular o segredo $S$, Eva precisaria encontrar $k_A$ ou $k_B$ a partir de $G$, o que é exatamente o \textbf{ECDLP}. Como o ECDLP é um problema difícil, o sistema é seguro.

% Parte escrita pelo JV

\section{Análise Comparativa: RSA versus Curvas Elípticas}
Após a apresentação dos sistemas RSA e da Criptografia de Curvas Elípticas (ECC), esta seção estabelece uma análise comparativa, destacando a eficiência e a robustez teórica de cada abordagem. A ECC representa uma evolução significativa em relação aos sistemas tradicionais como o RSA, baseando sua segurança em problemas matemáticos de complexidade distinta.

\subsection{Fundamentação Matemática e Complexidade}
A principal diferença entre os dois sistemas reside na dificuldade computacional dos problemas matemáticos que protegem as chaves.

\subsubsection{RSA: Fatoração de Inteiros}
A segurança do RSA depende da dificuldade de fatorar um número grande $n$, produto de dois primos $p$ e $q$:
\begin{equation}
n = p \cdot q
\end{equation}
O algoritmo mais eficiente conhecido para quebrar o RSA é o \textit{General Number Field Sieve} (GNFS). Sua complexidade de tempo para um inteiro $n$ é \textbf{subexponencial}, descrita aproximadamente por:
\begin{equation}
 L_n\left[\frac{1}{3}, c\right] = \exp\left( c \cdot (\ln n)^{\frac{1}{3}} \cdot (\ln \ln n)^{\frac{2}{3}} \right)
\end{equation}
Isso significa que, embora a dificuldade aumente com o tamanho da chave, ela não cresce de forma abrupta o suficiente. Por isso, sistemas RSA precisam de chaves gigantescas para compensar a eficiência desse algoritmo de ataque.

\subsubsection{ECC: Logaritmo Discreto em Curvas Elípticas (ECDLP)}
A segurança da ECC baseia-se em encontrar um inteiro $k$ (o logaritmo discreto), dado os pontos $P$ (ponto base) e $Q$ (chave pública) na curva, tal que:
\begin{equation}
 Q = k \cdot P
\end{equation}
Ao contrário da fatoração, não existem algoritmos subexponenciais conhecidos para resolver o ECDLP em curvas elípticas bem escolhidas. O melhor ataque genérico conhecido é o algoritmo \textit{Pollard's rho}, cuja complexidade é puramente \textbf{exponencial} (baseada na raiz quadrada do tamanho do grupo $n$):
\begin{equation}
\mathcal{O}(\sqrt{n}) = \mathcal{O}(2^{\frac{k}{2}})
\end{equation}
Onde $k$ é o número de bits da chave.

\subsection{Eficiência no Tamanho das Chaves}
Matematicamente, a ausência de um "atalho" subexponencial para a ECC (como o GNFS é para o RSA) explica a densidade superior de sua segurança. Enquanto na RSA aumentar a chave adiciona segurança de forma decrescente, na ECC, cada bit adicionado à chave dobra a dificuldade do ataque (curva de segurança íngreme). Isso resulta em uma segurança equivalente com chaves drasticamente menores, conforme ilustrado na Tabela \ref{tab:comparativo}.

\begin{table}[h]
\centering
\caption{Comparativo de Tamanho de Chaves (Bits) para Segurança Equivalente}
\vspace{0.3cm}
\begin{tabular}{ccc}
\toprule
\textbf{Nível de Segurança} & \textbf{Chave RSA} & \textbf{Chave ECC} \\
(Força Bruta Simétrica) & (GNFS Subexponencial) & (Pollard Exponencial) \\
\midrule
80 & 1024 & 160-192 \\
112 & 2048 & 224-256 \\
128 & 3072 & 256-384 \\
192 & 7680 & 384-512 \\
256 & 15360 & 512+ \\
\bottomrule
\end{tabular}
\label{tab:comparativo}
\end{table}

\subsection{Vantagens da ECC}
A superioridade da ECC reside na sua eficiência e no seu poder teórico. A robustez do problema do Logaritmo Discreto em Curvas Elípticas permite que a ECC ofereça, com 256 bits, uma segurança que o RSA só alcançaria com chaves na faixa de 15.360 bits. Esta redução massiva no tamanho das chaves se traduz em:
\begin{itemize}
\item \textbf{Processamento mais Rápido:} Assinaturas digitais e trocas de chaves mais rápidas.
\item \textbf{Menor Consumo de Energia:} Essencial para dispositivos móveis, IoT (Internet das Coisas) e ambientes com recursos limitados.
\item \textbf{Tamanho de Mensagem Reduzido:} Certificados digitais e cabeçalhos de TLS/SSL muito menores.
\end{itemize}
Isso torna a ECC o padrão ideal e preferido para a criptografia moderna.

\end{document}